\section{A fixed-temperature oracle for noisy regression}
\label{sec:noisy-regression}
We now prove Theorem~\ref{thm:fast-oracle}. The central quantity is the excess risk above the regression function, which removes the irreducible noise from the complexity penalty.

Write $f_0(x)=\E[Y\mid X=x]$, $\xi=Y-f_0(X)$, and
\[
 R_0=\E\xi^2,\qquad
 A(f)=\E(f(X)-f_0(X))^2=R(f)-R_0.
\]
The assumption is $|f_0|\le1$ and
\begin{equation}
 \E[e^{t\xi}\mid X]\le e^{\sigma^2t^2/2}
 \quad\text{for every }t\in\R.
 \label{eq:conditional-subgaussian}
\end{equation}
No Gaussian assumption is needed. Bounded labels in $[-1,1]$ satisfy this with $\sigma^2=1$ by conditional Hoeffding. Gaussian label noise of variance at most $\sigma^2$ is another example.

Let $s=1+\sigma^2$ and $\beta=\eta/(32s)$, where $0<\eta\le1$. Every network output $f_w$ remains bounded, although the labels need not be. The posterior is still~\eqref{eq:posterior}; its Gaussian observation variance is the constant $16s/\eta$, independent of sample size and learner width.

\begin{lemma}[A relative-loss PAC--Bayes inequality]\label{lem:relative-pac}
With probability at least $1-\delta/2$, simultaneously for every probability measure $\rho$ on weights,
\begin{equation}
 c_\beta^{\mathrm{rel}}\E_\rho A(f_w)
 \le\E_\rho[\whR_S(f_w)-\whR_S(f_0)]
       +\frac{\KL(\rho\|\mu_N)+\log(2/\delta)}{\beta n},
 \label{eq:relative-pac}
\end{equation}
where
\[
 c_\beta^{\mathrm{rel}}=\frac{1-e^{-4\beta(1-2\beta\sigma^2)}}{4\beta}
       \ge1-4s\beta=1-\eta/8.
\]
\end{lemma}
\begin{proof}
Fix $f=f_w$ and put $e=f-f_0$. The difference of squared losses is $e(X)^2-2\xi e(X)$. Conditional sub-Gaussianity gives
\[
 \E[e^{-\beta\{(f(X)-Y)^2-(f_0(X)-Y)^2\}}\mid X]
 \le e^{-\beta(1-2\beta\sigma^2)e(X)^2}.
\]
Since $0\le e^2\le4$, convexity on that interval yields
\[
 \E e^{-\beta\{(f(X)-Y)^2-(f_0(X)-Y)^2\}}
 \le e^{-\beta c_\beta^{\mathrm{rel}} A(f)}.
\]
The lower bound on $c_\beta^{\mathrm{rel}}$ follows from $1-e^{-u}\ge u-u^2/2$ and the chosen range of $\beta$. Independence, integration against the prior, Markov's inequality and the entropy variational inequality prove~\eqref{eq:relative-pac}. The individual squared losses have finite expectation under~\eqref{eq:conditional-subgaussian}; their difference has the exponential moment just established.
\end{proof}

The comparison measure may itself be averaged before applying concentration. This is useful because the perturbation lemma is naturally a statement about its expected risk.

\begin{lemma}[Concentration for a fixed comparison law]\label{lem:comparison-concentration}
Let $\nu$ be any fixed probability measure on bounded predictions, independent of the data, and let $a=\E_\nu A(f_w)$. With probability at least $1-e^{-\ell}$,
\begin{equation}
 \E_\nu[\whR_S(f_w)-\whR_S(f_0)]
 \le a+C\sqrt{sa\ell/n}+Cs\ell/n.
 \label{eq:comparison-bernstein}
\end{equation}
\end{lemma}
\begin{proof}
Set
\[
 v(x)=\E_\nu(f_w(x)-f_0(x))^2,\qquad
 b(x)=\E_\nu(f_w(x)-f_0(x)).
\]
Then $0\le v\le4$, $b^2\le v$, and one comparison-averaged excess loss is $L=v(X)-2\xi b(X)$, of mean $a$. For $0\le t\le1/(16s)$, conditional sub-Gaussianity gives
\[
 \E e^{t(L-a)}
 \le e^{-ta}\E e^{(t+2\sigma^2t^2)v(X)}
 \le\exp\left\{a\left[\frac{e^{4(t+2\sigma^2t^2)}-1}{4}-t\right]\right\}
 \le e^{6sa t^2}.
\]
For the last inequality, put $u=t+2\sigma^2t^2\le9t/8$, use $4u\le1$ and $e^{4u}-1\le4u+16u^2$. Chernoff's bound for the independent observations then gives~\eqref{eq:comparison-bernstein}, with an absolute constant. Equivalently, optimize over $t\le1/(16s)$ to obtain a square-root term and a linear term.
\end{proof}

\begin{proof}[Proof of Theorem~\ref{thm:fast-oracle}]
Fix a deterministic risk-minimizing $g\in\calG_{m,B}$ and write $A_*=A(g)$. For a precision $0<r\le1$, Lemma~\ref{lem:comparison}, applied pointwise with label $y=f_0(x)$, supplies a data-independent comparison law with
\[
 a=\E_\nu A(f_w)\le A_*+r,
 \qquad \KL(\nu\|\mu_N)\le K_{m,B}r^{-q}.
\]
This step does not require the observed labels to be bounded.

At the Gibbs posterior, the right side of~\eqref{eq:relative-pac} is
$F_{S,\beta}-\whR_S(f_0)+\log(2/\delta)/(\beta n)$.
The free-energy variational formula bounds it by
\[
 \E_\nu[\whR_S(f_w)-\whR_S(f_0)]
       +\frac{K_{m,B}r^{-q}+\log(2/\delta)}{\beta n}.
\]
Use Lemma~\ref{lem:comparison-concentration} with $\ell=\log(2/\delta)$ and Young's inequality to obtain, on the intersection of the two events,
\[
 c_\beta^{\mathrm{rel}}\E_{\Pi_{S,\beta}}A(f_w)
 \le(1+\eta/8)(A_*+r)
       +\frac{Cs\ell}{\eta n}
       +\frac{32sK_{m,B}r^{-q}}{\eta n}.
\]
Since $c_\beta^{\mathrm{rel}}\ge1-\eta/8$, the coefficient of $A_*$ after division is at most $1+\eta$. Let
\[
 u=\left(\frac{s\{K_{m,B}+\log(2/\delta)\}}{\eta n}\right)^{1/(q+1)}.
\]
If $u\le1$, take $r=u$; every remaining term is at most $C_q u$. If $u>1$, the deterministic bound $A(f_w)\le4$ gives the result after increasing $C_q$. Finally $R(f_w)=R_0+A(f_w)$, and Jensen's inequality transfers the posterior-average bound to the ordinary final posterior mean.
\end{proof}

The unknown regression function appears only inside the proof. The learner uses neither $f_0$, $R_0$, nor the comparator. The prescribed likelihood variance depends only on the chosen approximation factor $1+\eta$ and a known conditional sub-Gaussian proxy $\sigma^2$; an upper bound on that proxy is sufficient.
